\subsection{Коровы в стойла}

\subsubsection{Описание решения}

Применим стратегию бинарного поиска по ответу. Мы знаем, что ответ лежит в диапазоне от нуля до разности первого и последнего элементов исходного отсортированного массива. Предполагаем, что ответ лежит посередине данного отрезка. Если он подходит по условию, то двигаем левую границу к середине. Если нет --- двигаем правую границу.

Проверить корректность предположений просто: достаточно пробежаться циклом по всему массиву и пытаться расставить коров хотя бы на предполагаемом расстоянии друг от друга. Если коров получилось не меньше, чем в условии --- предположение верно, иначе --- нет.

В ответ пойдёт максимальное из предположений искомого расстояния.

\subsubsection{Асимптотическая сложность}

Временная сложность решения --- линейно-логарифмическая \( \mathcal{O}(N\cdot\log N) \). Бинарный поиск происходит за логарифмическое число шагов, причём каждый шаг выполняется за линейное время из-за обхода всего массива.

\subsubsection{Листинг кода}

\begin{lstlisting}
#include <iostream>
#include <vector>

int count_cows(std::vector<long>& arr, int n, long dist) {
  int count = 1;
  long last = arr[0];
  for (int i = 1; i < n; i++) {
    if (std::abs(last - arr[i]) >= dist) {
      last = arr[i];
      count++;
    }
  }
  return count;
}

long bin_search(std::vector<long>& arr, int n, int k) {
  long l = 0, r = arr[n - 1] - arr[0];
  while (l + 1 < r) {
    long mid = l + (r - l) / 2;
    if (count_cows(arr, n, mid) < k) {
      r = mid - 1;
    } else {
      l = mid;
    }
  }
  if (count_cows(arr, n, r) >= k) {
    return r;
  }
  return l;
}

int main() {
  int N, K;
  std::cin >> N >> K;
  std::vector<long> arr;
  for (int i = 0; i < N; i++) {
    long item;
    std::cin >> item;
    arr.push_back(item);
  }
  std::cout << bin_search(arr, N, K);
  return 0;
}
\end{lstlisting}